\section{Task 1: FER2013 baseline}
\label{sec:task1}
% Numbers are taken from the saved notebook outputs (notebooks/01_task1_fer2013_models.ipynb) and figures/*.csv.

\textbf{Features.} Three feature types were extracted, each compared with three classifiers ($3\times3=9$ models):
\begin{itemize}
\item \textbf{HOG (900 dimensions):} histogram of oriented gradients on the $48\times48$ image, 9 orientations, $8\times8$ pixel cells, $2\times2$ blocks with L2-Hys normalization. HOG encodes local edge orientation (mouth corners, brows, eye shape, wrinkles) and the block normalization gives some robustness to illumination.
\item \textbf{Facial landmarks (136 dimensions):} the 68 dlib landmarks, centred and scaled to $[-1, 1]$. The whole $48\times48$ image is passed to dlib as the face box, so the same procedure can be reused in the real-time demo.
\item \textbf{PCA (100 dimensions):} raw pixels scaled to $[0, 1]$, projected with whitened PCA. PCA is fitted inside the scikit-learn pipeline on the training data only, which avoids leakage into the validation and test sets.
\end{itemize}

\textbf{Classifiers.} Multinomial Logistic Regression ($C = 1$, standardized inputs), an RBF-kernel SVM ($C = 1$, $\gamma =$ ``scale'', standardized inputs) and a Random Forest (300 trees, sqrt features, no scaling). All use a fixed random seed (42).

\textbf{Results on the validation set.} Table~\ref{tab:fer-val} in Appendix~\ref{app:tables} shows accuracy, balanced accuracy and macro-F1 for all nine combinations. Macro-F1 was the selection criterion, because it weights all seven classes equally and is therefore not dominated by Happy.

\textbf{Best classifier.} The RBF-SVM is the best classifier on two of the three feature types (HOG: 0.505 macro-F1; PCA: 0.430) and the best overall. Its non-linear kernel can model the non-linear relation between local gradient patterns and expression, which the linear Logistic Regression cannot: Logistic Regression is the weakest or near-weakest classifier for every feature type (0.287 to 0.395 macro-F1). The Random Forest is competitive only on landmarks, where it is the best classifier (0.474) and clearly beats the SVM (0.401). A plausible but untested reason for the landmark result is given in Appendix~\ref{app:feature-discussion}.The SVM is also by far the most expensive model to train (438 s on HOG, against 36 s for the HOG forest and 23 s for Logistic Regression).

\textbf{Best features.} HOG is the best representation for Logistic Regression and the SVM (for example 0.505 against 0.430 for PCA and 0.401 for landmarks with the SVM). For the Random Forest, landmarks are the best representation (0.474 against 0.440 for HOG). We attribute the HOG result to its local edge information, and the reasons for each ranking, including the sensitivity of landmarks and PCA to misalignment, are discussed in Appendix~\ref{app:feature-discussion}.

\textbf{Hyperparameter tuning and final test result.} Each classifier was tuned on the feature set where it did best on validation, and tuning raised the SVM from 0.505 to 0.562 validation macro-F1. The search grids and the tuned validation scores for every classifier are given in Appendix~\ref{app:tuning}. Each tuned model was then evaluated once on the private test split (Table~\ref{tab:fer-test}).

\begin{table}[h]
\centering\small
\caption{Final test performance of each tuned classifier on the PrivateTest split, with 95\% bootstrap intervals for macro-F1 (1,000 resamples). The confusion matrices are in Appendix~\ref{app:cm}.}
\label{tab:fer-test}
\begin{tabular}{@{}llcccc@{}}\toprule
Features & Classifier & Accuracy & Bal. acc. & Macro-F1 & 95\% interval \\\midrule
HOG & SVM ($C = 10$) & 0.575 & 0.550 & 0.568 & [0.544, 0.590] \\
Landmarks & Random Forest & 0.513 & 0.463 & 0.482 & [0.457, 0.507] \\
HOG & Logistic Regression & 0.438 & 0.376 & 0.378 & [0.355, 0.400] \\
\bottomrule\end{tabular}
\end{table}

The test score of the SVM is higher than its validation score (0.568 against 0.562 tuned, and 0.505 default). The paired bootstrap difference between the SVM and the landmark Random Forest is $+0.086$ with interval $[+0.062, +0.111]$, so the SVM's lead on the test split is unlikely to be noise. The gap between accuracy and balanced accuracy on validation (0.542 against 0.482) is caused mainly by the low recall of Disgust (25\%).

\textbf{Which emotions are hardest to distinguish?} Averaging per-class F1 over the nine validation models gives the following ordering from hardest to easiest: Fear (0.26), Disgust (0.27), Angry (0.32), Sad (0.38), Neutral (0.43), Surprise (0.62), Happy (0.64). For the SVM on the test set, F1 is 0.75 for Happy and 0.73 for Surprise, and 0.43 to 0.46 for Angry, Sad and Fear.

\begin{itemize}
\item Happy and Surprise are relatively easier to classify. Both have distinctive, large-scale configurations: a raised, widened mouth with visible teeth for Happy, and wide-open eyes and mouth for Surprise. Happy additionally has the most training data.

\item Fear, Angry and Sad are harder and are confused with each other and with Neutral. The most frequent errors of the validation SVM are listed in Appendix~\ref{app:feature-discussion}; the test set shows the same pattern (Appendix~\ref{app:cm}). These classes are expressed through subtle, low-intensity changes around the brows and eyes, which are hard to resolve at $48\times48$ pixels. Fear also shares features with Surprise (wide eyes, open mouth), and Sad and Neutral faces are visually close, so Neutral is a frequent ``sink'' class.

\item Disgust is rare (436 training images), so the classifier rarely predicts it. On validation the SVM has precision 1.00 but recall 0.25: when it predicts Disgust it is right, but it misses most Disgust faces. On the test set recall is 0.51, but with only 55 test images this estimate is noisy. Class balancing was tested in Task 5 and did not help.
\end{itemize}

\textbf{Evidence from error analysis.} Illumination and contrast have a small effect on error, around five percentage points: the SVM error rate is about 0.46 to 0.49 in the three darkest brightness quintiles and about 0.42 to 0.44 in the two brightest (Figure~\ref{fig:err-light}, Appendix~\ref{app:figs}). Recurring error types, including ambiguous or questionable labels, occlusion, non-frontal pose and overlaid text or watermarks, are illustrated in Figure~\ref{fig:err-gallery} (Appendix~\ref{app:figs}). Both figures are qualitative evidence rather than counts.

\textbf{Duplicates.} FER2013 contains exact copies of images across splits, and some copies carry different labels, which is direct evidence of label noise. Duplicates inflate the test score: on the 3,301 non-duplicated test images the SVM reaches 0.505 macro-F1, against 0.568 on the full test split, so 0.505 is the fairer estimate. The duplicate breakdown and the validation effect of removing duplicates are given in Appendix~\ref{app:dup}. We treat cleaning the training set as a data curation step in Task 5.

\textbf{Summary and limitations.} The tuned HOG + RBF-SVM is the best FER2013 baseline (test: 0.575 accuracy, 0.550 balanced accuracy, 0.568 macro-F1; 0.505 macro-F1 without duplicates). HOG is the strongest feature type for Logistic Regression and the SVM, and landmarks are the strongest for the Random Forest; the SVM is the strongest classifier on HOG and PCA. Happy and Surprise are recognized reliably, while Fear, Angry, Sad and Disgust are not. Prediction times for the demo are given in Appendix~\ref{app:feature-discussion}.Limitations: (i) the SVM was tuned on a 10,000-image subsample; (ii) the final test was run once per model, with no further selection; (iii) no class balancing was used, which disadvantages Disgust; (iv) the landmark features depend on a detector applied to $48\times48$ images with the whole image as the face box.
